\scanpage{5}\mspage{4}
D'autre part, il y a deux \inserted{\uline{autres}} applications
\[
\left\{\begin{aligned}
\text{a)}\quad &\unsure{\mathcal{F}|\mathcal{E}}\longrightarrow\unsure{\mathcal{P}|\mathscr{P}}(\deleted{\unsure{\mathcal{F}}}\mathcal{M})&
 F&\longmapsto\operatorname{Omb}_{\unsure{\mathcal{M}|m}}(F)\\
\text{b)}\quad &\mathcal{F}\hookrightarrow\mathcal{P}\mathcal{P}(\deleted{\unsure{\mathcal{F}}}\mathcal{M})&
 \text{et }\unsure{=}\quad\mathcal{F}&\longrightarrow\operatorname{Fig}(\mathcal{M}\unsure{]|)}\subset\mathcal{P}\mathcal{P}\unsure{\mathcal{M}|(\mathcal{M})}\\
&& F&\mapsto\deleted{\unsure{O}}\operatorname{Multomb}_{\mathcal{M}}(F),
\end{aligned}\right.\tag*{(1.\deleted{\unsure{4}}6)}
\]
où $\operatorname{Omb}_{\mathcal{M}}(F)$ s'appelle l'\uline{\unsure{ogombre|ombre}} \uline{de la} figure $F$ sur $\mathcal{M}$, ou la \uline{grande} \uline{ombre} \uline{de} $\mathcal{F}$. \deleted{\unsure{La}}Une multistrate \inserted{$X$} qui est dans la grande ombre d'une figure $F$ est dite ``\uline{raffiner}'' \uline{la figure} $F$ (c'est cette fois déduit d'une relation d'ordre $\ll$ sur $\mathbb{F}$, le raffinement, par la condition $X\ll F$). Il n'est pas clair que cette application $\mathcal{F}\unsure{\relbar|\to}\mathcal{P}(\mathcal{M})$ soit injective, ie. que $\mathcal{F}$ s'identifie à un \unsure{ens} de parties de $\mathcal{M}$, via cette application. Par contre il en sera ainsi de (1.6) b), qui est déduit de a) via \deleted{\unsure{$F$}}la \unsure{formation|fonction} de $\widetilde F$, comme (1.3 b) est déduit de (1.3. a):
\[
\operatorname{Multomb}(F)=\{X\in\widetilde F\mid\operatorname{Omb}X\},\tag*{(1.\deleted{\unsure{5|6}}7)\quad\deleted{\unsure{(4.5)|(1.5)}}}
\]
alors que $\operatorname{Omb}(X)$ s'explicite par
\[
\unsure{\operatorname{Multomb}|\operatorname{MultOmb}}(X)=\{Y\in\deleted{\unsure{\mathcal{F}}}\mathcal{M}\mid Y\ll X\},\tag*{(1.\deleted{\unsure{6}}8)}
\]
analogue à (1.5), où cette fois on a utilisé la relation \inserted{d'ordre} $\ll$ ald $\leq$ -- \deleted{\unsure{C'est là}}
